\clearpage
\section{Spinning Sphere}
\subsection{Surface Current density}
velocity on a rotating sphere is 
\begin{equation}
  \vb{v} = \vb{\omega}\times\vb{r}
\end{equation}
but in spherical
\begin{equation}
  \vb{\omega}\times\vb{r}=\omega R \sin\theta\vu{\varphi}
\end{equation}
so 
\begin{equation}
  \vb{K} = \sigma\omega R \sin\theta\vu{\varphi}
\end{equation}
and the volumetric density is 
\begin{equation}
  \vb{j} = \sigma\omega R \sin\theta\delta\ins{r-R}\vu{\varphi}
\end{equation}
\subsection{magnetic dipole}
we know that 
\begin{equation}
  \vb{m} = \frac{1}{2c}\int \dd{x}^3 \vb{r}\times j 
\end{equation}
since the volumetric integral will become a surface integral using the delta, we can say 
\begin{equation}
  \vb{m} = \frac{1}{2c}\oint \vb{r}\times \vb{K}\dd{A}
\end{equation}
but 
\begin{equation}
  \vb{r}\times\vb{K}=\sigma\omega R^2\sin\theta \ins{\vu{r}\times\vu{\varphi}}=-\sigma\omega R^2\sin\theta\vu{\theta}
\end{equation}
but we know the moment must point in the \(\vu{z}\) direction here, so 
\begin{equation}
  \ins{\vb{r}\times\vb{K}}\cdot\vu{z}=\sigma\omega R^2\sin^2\theta
\end{equation}
thus 
\begin{align}
  m&=\frac{\sigma\omega R^2}{2c}\oint\sin^2\theta\dd{A}\\
  &= \frac{\sigma\omega \pi R^4}{c}\int_{0}^{\pi}\sin^3\theta\dd{\theta}\\
  &= \frac{\sigma\omega \pi R^4}{c}\cdot\frac{4}{3}\\
  &= \frac{4\sigma\omega \pi R^4}{3c}
\end{align}
so 
\begin{equation}
  \vb{m} = \frac{4\pi\sigma\omega R^4}{3c}\vu{z}
\end{equation}
\subsection{fields inside and outside}
in TA 10 we saw that 
\begin{equation}
  \vb{A}\ins{\vb{r}}=\frac{4\pi\sigma\omega R^3}{3c}\frac{r_<}{r_>^2}\sin\theta\vu{\varphi}
\end{equation}
which we can rewrite using the magnetic dipole 
\begin{equation}
  \vb{A}\ins{\vb{r}} = \begin{cases}
    \frac{4\pi\sigma\omega R}{3c}r\sin\theta\vu{\varphi} & r<R\\
    \frac{4\pi\sigma\omega R^4}{3cr^2}\sin\theta\vu{\varphi} & r>R
  \end{cases}
\end{equation}
or
\begin{equation}
  \vb{A}\ins{\vb{r}} = \begin{cases}
    \frac{mr}{R^3}\sin\theta\vu{\varphi} & r<R\\
    \frac{m}{r^2}\sin\theta\vu{\varphi} & r>R
  \end{cases}
\end{equation}
The magnetic field is the rotor of the vector potential, so we know that, using the spherical rotor
\begin{equation}
  B_r = \frac{1}{r\sin\theta}\partial_\theta\ins{\sin\theta A_\varphi},\qquad B_\theta = -\frac{1}{r}\partial_r\ins{r A_\varphi},\qquad B_\varphi=0
\end{equation}
Looking inside the sphere, 
\begin{equation}
  B_r = \frac{1}{r\sin\theta}\partial_\theta\ins{m\frac{r}{R^3}\sin^2\theta}=\frac{2m}{R^3}\cos\theta
\end{equation}
and 
\begin{equation}
  B_\theta = -\frac{1}{r}\partial_r\ins{\frac{mr^2}{R^3}\sin\theta\vu{\varphi}}=-\frac{2m}{R^3}\sin\theta
\end{equation}
so 
\begin{equation}
  \vb{B}=\frac{2m}{R^3}\ins{\cos\theta\vu{r}-\sin\theta\vu{\theta}}
\end{equation}
but using the identity
\begin{equation}
  \cos\theta\vu{r}-\sin\theta\vu{\theta}=\vu{z}
\end{equation}
\begin{equation}
  \vb{B}=\frac{8\pi\sigma\omega R}{3c}\vu{z}
\end{equation}
and outside the sphere 
\begin{equation}
  B_r = \frac{1}{r\sin\theta}\partial_\theta\ins{\frac{m}{r^2}\sin^2\theta}=2\frac{m}{r^3}\cos\theta
\end{equation}
and 
\begin{equation}
  B_\theta = -\frac{1}{r}\partial_r\ins{\frac{m}{r}\sin\theta}=\frac{m}{r^3}\sin\theta
\end{equation}
so 
\begin{equation}
  \vb{B}=\frac{m}{r^3}\ins{2\cos\theta\vu{r}+\sin\theta\vu{\theta}}
\end{equation}
